% Part: incompleteness
% Chapter: representability-in-q
% Section: basic-representable

\documentclass[../../../include/open-logic-section]{subfiles}

\begin{document}

\olfileid{inc}{req}{bre}
\olsection{મૂળભૂત વિધેયો~$\Th{Q}$માં નિરૂપ્ય છે}

પહેલાં બધા મૂળભૂત વિધેયો
$\Th{Q}$માં નિરૂપ્ય છે તે બતાવવું
પડશે. છેવટે દરેક $k$-ચલવાળા
મૂળભૂત વિધેય $f(x_0,\dots,x_{k-1})$
માટે તેનું નિરૂપણ કરતું
!!a{formula}
$!A_f(x_0,\dots,x_{k-1},y)$
કેવી રીતે આપવું તે બતાવવાનું છે.

શૂન્ય વિધેય, ઉત્તરગામી વિધેય,
સરવાળો, ગુણાકાર, સમતા માટેનું
લાક્ષણિક વિધેય અને પ્રક્ષેપણોનું
નિરૂપણ કરી શકીશું. દરેક કિસ્સામાં
યોગ્ય નિરૂપક સૂત્ર સીધું છે:
દાખલા તરીકે, શૂન્ય વિધેયને
$y = \Obj 0$ સૂત્ર, ઉત્તરગામીને
$x_0' = y$ !!{formula}, અને
સરવાળાને $(x_0 + x_1) = y$
!!{formula} નિરૂપે છે. કામ એ
બતાવવાનું છે કે સંબંધિત
!!{sentence}s~$\Th{Q}$માં
સાબિત થાય છે. દાખલા તરીકે,
ઉપરનું !!{formula} સરવાળાને
નિરૂપે છે એમ કહેવા માટે દરેક
પ્રાકૃતિક સંખ્યાની જોડ $m$ અને
$n$ માટે $\Th{Q}$માં આ બંને
સાબિત થાય છે તે બતાવવું પડે:
\begin{align*}
& \eq[\num n + \num m][\num {n+m}] \text{અને}\\
& \lforall[y][(\eq[(\num n + \num m)][y] \lif \eq[y][\num{n+m}])].
\end{align*}

\begin{prop}
\ollabel{prop:rep-zero}
શૂન્ય વિધેય $\Zero(x) = 0$ને
$!A_{\Zero}(x,y) \ident \eq[y][\Obj 0]$
સૂત્ર~$\Th{Q}$માં નિરૂપે છે.
\end{prop}

\begin{prop}
\ollabel{prop:rep-succ}
ઉત્તરગામી વિધેય $\Succ(x) = x+1$ને
$!A_{\Succ}(x,y) \ident \eq[y][x']$
સૂત્ર~$\Th{Q}$માં નિરૂપે છે.
\end{prop}

\begin{prop}
\ollabel{prop:rep-proj}
પ્રક્ષેપણ વિધેય
$\Proj{n}{i}(x_0, \dots, x_{n-1}) = x_i$ને
$\Th{Q}$માં આ સૂત્ર નિરૂપે છે:
\[!A_{\Proj{n}{i}}(x_0, \dots, x_{n-1}, y) \ident \eq[y][x_i].\]
\end{prop}

\begin{prob}
$\eq[y][\Obj 0]$, $\eq[y][x']$
અને $\eq[y][x_i]$ અનુક્રમે
$\Zero$, $\Succ$ અને
$\Proj{n}{i}$ને નિરૂપે છે
તે સાબિત કરો.
\end{prob}

\begin{prop}
\ollabel{prop:rep-id}
સમતા~$=$નું લાક્ષણિક વિધેય
\[
\Char{=}(x_0, x_1) =
\begin{cases}
  1 & \text{જો } x_0 =x_1\\
  0 & \text{અન્યથા}
\end{cases}
\]
ને~$\Th{Q}$માં આ સૂત્ર નિરૂપે છે:
\[
  !A_{\Char{=}}(x_0, x_1, y) \ident (\eq[x_0][x_1] \land \eq[y][\num{1}]) \lor (\eq/[x_0][x_1] \land
\eq[y][\num{0}]).
\]
\end{prop}

સાબિતી માટે આ લેમ્મા જરૂરી છે.

\begin{lem}
\ollabel{lem:q-proves-neq}
પ્રાકૃતિક સંખ્યાઓ $n$ અને $m$
માટે, જો $n \neq m$ હોય, તો
$\Th{Q} \Proves \eq/[\num n][\num m]$.
\end{lem}

\begin{proof}
$n$ પર આગમન વડે સાબિત કરો કે
દરેક $m$ માટે, $n \neq m$ હોય
તો $Q \Proves \eq/[\num n][\num m]$.

પાયાના કિસ્સામાં $n = 0$.
$m$ એ $0$ની બરાબર ન હોય, તો
કોઈ પ્રાકૃતિક સંખ્યા $k$ માટે
$m = k + 1$. આપણી પાસે
$\lforall[x][\eq/[0][x']]$
કહેતું સ્વયંસિદ્ધ વાક્ય છે.
પરિમાણકના સ્વયંસિદ્ધ વાક્યથી
$x$ને $\num k$ વડે બદલીને
$\eq/[0][\num k']$ તારવી
શકીએ. પરંતુ $\num k'$ એ જ
$\num m$ છે.

આગમનના પગથિયામાં $n$ માટે
દાવો સાચો છે એમ માનીને $n+1$
વિચારીએ. $m$ કોઈપણ પ્રાકૃતિક
સંખ્યા લો. બે શક્યતાઓ છે:
$m = 0$, અથવા કોઈ $k$ માટે
$m = k+1$. પહેલો કિસ્સો ઉપર
પ્રમાણે ઉકેલાય છે. બીજા કિસ્સામાં
માનો કે $n+1 \neq k+1$.
તો $n \neq k$. $n$ માટેની
આગમન-ધારણાથી
$\Th{Q} \Proves \eq/[\num n][\num k]$.
આપણી પાસે આ સ્વયંસિદ્ધ વાક્ય છે:
$\lforall[x][\lforall[y][\eq[x'][y'] \lif \eq[x][y]]]$.
પરિમાણકના સ્વયંસિદ્ધ વાક્યથી
$\eq[\num n'][\num k'] \lif \eq[\num n][\num
  k]$ મળે છે. વિધાનાત્મક તર્કથી
$\Th{Q}$માં
$\eq/[\num n][\num k] \lif \eq/[\num n'][\num k']$
તારવી શકીએ. મોડસ પોનેન્સથી
$\eq/[\num n'][\num k']$ મળે,
અને એ જ જોઈએ છે, કારણ કે
$\num k'$ એ $\num m$ છે.
\end{proof}

\begin{explain}
લેમ્માનો દાવો બહુ મોટો નથી:
તેનો સાર એટલો છે કે જુદાં
અંકનિરૂપણો જુદા પદાર્થો દર્શાવે
છે તે $\Th{Q}$ સાબિત કરી શકે છે.
દાખલા તરીકે, $\Th{Q}$ આ સાબિત
કરે છે: $0'' \neq 0'''$.
પરંતુ આ સામાન્ય રીતે સાચું છે
તે બતાવવા થોડી કાળજી જોઈએ.
આગમન વાપરીએ છીએ, પણ એ
$\Th{Q}$ની \emph{બહાર}નું
આગમન છે, એ પણ ધ્યાન રાખો.
\end{explain}

\begin{proof}[\olref{prop:rep-id}ની સાબિતી]
$n = m$ હોય, તો $\num{n}$ અને
$\num{m}$ એક જ પદ છે અને
$\Char{=}(n, m) = 1$.
પરંતુ $\Th{Q} \Proves
(\eq[\num{n}][\num{m}] \land
\eq[\num{1}][\num{1}])$,
તેથી $!A_=(\num{n}, \num{m},
\num{1})$ પણ સાબિત થાય છે.
$n \neq m$ હોય, તો
$\Char=(n, m) = 0$.
\olref{lem:q-proves-neq} પ્રમાણે
$\Th{Q} \Proves \eq/[\num{n}][\num{m}]$;
તેથી
$(\eq/[\num{n}][\num{m}] \land \Obj 0 = \Obj 0)$
પણ સાબિત થાય છે. આમ
$\Th{Q} \Proves !A_=(\num{n},
\num{m}, \num{0})$.

બીજા ભાગમાં પણ બે કિસ્સા છે.
$n = m$ હોય, તો
$\Th{Q} \Proves \lforall[y][(!A_=(\num{n}, \num{m}, y)
  \lif \eq[y][\num{1}])]$
બતાવવાનું છે. અનૌપચારિક રીતે
દલીલ કરીએ: માનો કે
$!A_=(\num{n}, \num{m}, y)$,
એટલે કે
\[
(\eq[\num{n}][\num{n}] \land \eq[y][\num{1}]) \lor
(\eq/[\num{n}][\num{n}] \land \eq[y][\num{0}])
\]
ડાબો વિકલ્પ તર્કથી
$\eq[y][\num{1}]$ આપે છે;
જમણો વિકલ્પ તર્કથી સાબિત થતા
$\eq[\num{n}][\num{n}]$ સાથે
વિરોધાભાસમાં છે.

બીજી તરફ, માનો કે $n \neq m$.
ત્યારે $!A_=(\num{n},
\num{m}, y)$ આ છે:
\[
(\eq[\num{n}][\num{m}] \land \eq[y][\num{1}]) \lor
(\eq/[\num{n}][\num{m}] \land \eq[y][\num{0}])
\]
અહીં ડાબો વિકલ્પ
$\eq/[\num{n}][\num{m}]$ સાથે
વિરોધાભાસમાં છે; આ વિધાન
$\Th{Q}$માં
\olref{lem:q-proves-neq}થી
સાબિત થાય છે. જમણા
વિકલ્પમાંથી $\eq[y][\num{0}]$
નીકળે છે.
\end{proof}

\begin{prop}
\ollabel{prop:rep-add}
સરવાળાનું વિધેય
$\Add(x_0, x_1) = x_0+x_1$ને
$\Th{Q}$માં આ સૂત્ર નિરૂપે છે:
\[
  !A_{\Add}(x_0, x_1, y) \ident \eq[y][(x_0 + x_1)].
\]
\end{prop}

\begin{lem}
\ollabel{lem:q-proves-add}
$\Th{Q} \Proves \eq[(\num{n} + \num{m})][\num{n+m}]$
\end{lem}

\begin{proof}
આ દાવો~$m$ પર આગમનથી
સાબિત કરીએ. $m = 0$ હોય, તો
$\Th{Q} \Proves
\eq[(\num{n} + \Obj 0)][\num{n}]$
બતાવવાનું છે. તે સ્વયંસિદ્ધ
વાક્ય~$!Q_4$થી મળે છે.
હવે $m$ માટે દાવો માનીને
$m+1$ માટે, એટલે કે
$\Th{Q} \Proves \eq[(\num{n} +
  \num{m+1})][\num{n+m+1}]$,
સાબિત કરીએ. ધ્યાન આપો કે
$\num{m+1}$ એ $\num{m}'$
અને $\num{n+m+1}$ એ
$\num{n+m}'$ જ છે.
સ્વયંસિદ્ધ વાક્ય $!Q_5$થી
$\Th{Q} \Proves
\eq[(\num{n} + \num{m}')][(\num{n}+\num{m})']$.
આગમન-ધારણાથી
$\Th{Q} \Proves
\eq[(\num{n} + \num{m})][\num{n+m}]$.
તેથી $\Th{Q} \Proves
\eq[(\num{n} + \num{m}')][\num{n+m}']$.
\end{proof}

\begin{proof}[\olref{prop:rep-add}ની સાબિતી]
$\Add$નું નિરૂપણ કરતું
!!{formula}~$!A_\Add(x_0, x_1, y)$
એ $\eq[y][(x_0 + x_1)]$ છે.
પહેલાં બતાવીએ કે
$\Add(n, m) = k$ હોય, તો
$\Th{Q} \Proves
!A_\Add(\num{n}, \num{m}, \num{k})$,
એટલે કે $\Th{Q} \Proves
\eq[\num{k}][(\num{n} + \num{m})]$.
પરંતુ $k = n + m$ હોવાથી
$\num{k}$ એ $\num{n+m}$ જ છે;
અને \olref{lem:q-proves-add}માં
$\Th{Q} \Proves \eq[(\num{n} +
  \num{m})][\num{n+m}]$
પહેલેથી બતાવ્યું છે.

$\Add(n, m) = k$ હોય, તો આ પણ
બતાવવું પડે:
\[
\Th{Q} \Proves \lforall[y][(!A_\Add(\num{n}, \num{m}, y) \lif
  \eq[y][\num{k}])].
\]
માનો કે
$\eq[(\num{n} + \num{m})][y]$ છે.
કારણ કે
\[
\Th{Q} \Proves \eq[(\num{n}+\num{m})][\num{n+m}],
\]
ડાબી બાજુને $\num{n+m}$ વડે
બદલતાં કોઈપણ~$y$ માટે
$\eq[\num{n+m}][y]$ મળે છે.
\end{proof}

\begin{prop}
\ollabel{prop:rep-mult}
ગુણાકારનું વિધેય
$\Mult(x_0, x_1) = x_0 \cdot x_1$ને
$\Th{Q}$માં આ સૂત્ર નિરૂપે છે:
\[
  !A_{\Mult}(x_0, x_1, y) \ident y = (x_0 \times x_1).
\]
\end{prop}

\begin{proof} અભ્યાસપ્રશ્ન. \end{proof}

\begin{lem}
\ollabel{lem:q-proves-mult}
$\Th{Q} \Proves \eq[(\num{n} \times \num{m})][\num{n \cdot m}]$
\end{lem}

\begin{proof} અભ્યાસપ્રશ્ન. \end{proof}

\begin{prob}
\olref[inc][req][bre]{lem:q-proves-mult}
સાબિત કરો.
\end{prob}

\begin{prob}
\olref[inc][req][bre]{lem:q-proves-mult}
વાપરીને
\olref[inc][req][bre]{prop:rep-mult}
સાબિત કરો.
\end{prob}

\begin{explain}
યાદ રાખો કે અંકગણિતની
ભાષામાં વિધેય-ચિહ્ન માટે
$\times$ અને સંખ્યાઓ પરની
સામાન્ય ગુણાકારક્રિયા માટે
$\cdot$ વાપરીએ છીએ.
તેથી સંખ્યાઓ દર્શાવતી
અભિવ્યક્તિઓની વચ્ચે
$\cdot$ આવી શકે છે (જેમ કે
$m \cdot n$માં), જ્યારે
$\times$ માત્ર અંકગણિતની
ભાષાનાં પદોની વચ્ચે આવે છે
(જેમ કે $(\num{m} \times
  \num{n})$માં). વાત વધુ
ગૂંચવાય છે, કારણ કે $+$
બંને માટે વપરાય છે:
!!{function} માટે પણ અને
સરવાળાની ક્રિયા માટે પણ.
તે પદોની વચ્ચે આવે—દાખલા
તરીકે $(\num{n} + \num{m})$માં—તો
તે અંકગણિતની ભાષાનું $2$-સ્થાની
!!{function} છે; અને સંખ્યાઓની
વચ્ચે આવે—દાખલા તરીકે
$n+m$માં—તો તે સરવાળાની
ક્રિયા છે. $\num{n+m}$નો
કિસ્સો પણ આમાં આવે છે:
તે સંખ્યા~$n+m$ને અનુરૂપ
પ્રમાણભૂત અંકનિરૂપણ છે.
\end{explain}

\end{document}
